\subsection*{\S. 23.\quad 达布方法。}
除了上一节所研究问题的解析解以外，现在还介绍第二种更富几何性的解法；这正是达布（\emph{同上引书}）给出的解法，他为此采用了 Combescure 变换（第20节）。

上一节给出了一个 \(\infty^1\) 球面族，以球心的轨迹曲线 C 和球面半径 \(R=R(s)\) 来确定它。与它一起，再考察任意另一个球面族：按照第20节的公式，通过 Combescure 变换把球心轨迹 C 变为另一条曲线 C\(_1\)：
\[
\diff s_1=f(s)\,\diff s,
\qquad
\rho_1=\rho f(s),
\qquad
T_1=T f(s),
\]
同时以同一比率改变球面半径，令
\[
R_1=R f(s);
\]
称第二个球面族由第一个球面族作 Combescure 变换而得。

立即可以看出，关于 \(\theta,\varphi\) 的方程 (48) 保持不变；所以，在积分给定球面族的正交轨迹问题时，也同时解决了所有由 Combescure 变换所得球面族的同一问题。这个简单结果的几何解释也很容易给出：两个对应球面上由相同 \(\theta,\varphi\) 值确定的两点，分别与各自球心相连，便给出两条平行半径。因此：\emph{由给定球面族的一条正交轨迹，可以推出 Combescure 变换球面族的一条正交轨迹：在第二个球面族的每个球面 S\(_1\) 上，标出点 M\(_1\)，它对应于原球面 S 与正交轨迹相交的点 M。}

所以可以看出，当一个 \(\infty^1\) 球面族作 Combescure 变换时，它们的正交轨迹也随之作 Combescure 变换。

达布方法所依据的基本说明如下：\emph{已知球面族的一条正交轨迹 C，就可以通过求积把这个球面族换成一个 Combescure 变换球面族，使其中所有球面都通过一个固定点。}

为证明这一点，对曲线 C 采用通常的记号，则
\[
x_0=x+R\alpha,
\qquad y_0=y+R\beta,
\qquad z_0=z+R\gamma
\]
是球心的坐标。由一个固定点，例如原点 O，引出一条线段 R\(_1\)（待定为 \(s\) 的函数），使它平行于 C 在 \((x,y,z)\) 处的切线；其端点坐标为
\[
x_1=R_1\alpha,
\qquad y_1=R_1\beta,
\qquad z_1=R_1\gamma
\]
为了使球面
\[
(X-x_1)^2+(Y-y_1)^2+(Z-z_1)^2=R_1^2
\]
描出 Combescure 变换球面族，必须把 R\(_1\) 确定为 \(s\) 的函数，使
\[
\frac{\diff x_1}{\diff s}:\frac{\diff y_1}{\diff s}:\frac{\diff z_1}{\diff s}
=
\frac{\diff x_0}{\diff s}:\frac{\diff y_0}{\diff s}:\frac{\diff z_0}{\diff s}.
\]
这些条件立即给出 R\(_1\) 的方程
\[
\frac{R'_1}{R_1}=\frac{1+R'}{R}
\]
因而积分得到
\[
R_1=C R e^{\int\frac{\diff s}{R}}\quad (C\hbox{ 为常数}).
\]
反过来，这样所得的球面族确实是原球面族的 Combescure 变换，因为立即可验证
\[
\frac{\diff s_1}{R_1}=\frac{\diff s_0}{R},
\]
其中 \(\diff s_0,\diff s_1\) 表示相应球心轨迹曲线的弧长元。

这样把球面族化为通过一个固定点以后，以该固定点为中心作向径互为倒数的反演，便把这些球面变成一个 \(\infty^1\) 平面族，并把球面的正交轨迹变成平面的正交轨迹；因为众所周知，这种反演保持角度。于是，一般问题最终化为第19节的特殊问题，并因此可以通过求积解决。
